\section{Introduction}\label{sec:intro}

Among Californians aged 25 to 64 in the 2022 American Community Survey (ACS), 88\% of positive wage and salary reports are multiples of \$1{,}000, 55\% of \$5{,}000 and 38\% of \$10{,}000; in the Current Population Survey (CPS) the shares are 86\%, 54\% and 35\%, and a further 15\% of reports are multiples of \$25{,}000. Such heaping perturbs a mean very little, but a kernel estimate of reported earnings has a local maximum at every round number, and a test of unimodality applied to the reports tests the questionnaire. In the simulations of Section~\ref{sec:sim}, the excess mass test of \citet{AmeijeirasAlonsoCRC2019} and the dip test of \citet{ChengHall1998} reject unimodality in every one of $16{,}800$ replications drawn from seven unimodal latent densities and reported through three rounding mechanisms, and hold their level with unrounded data wherever the latent density is unimodal on the scale on which they are applied.

Bimodality, polarization and a thinning middle of the income distribution are established topics in the study of inequality \citep{EstebanRay1994,Wolfson1994,Jenkins1995,Quah1996,CowellJenkinsLitchfield1996,DuclosEstebanRay2004,PittauZelliJohnson2010}, and all rest on the shape of a density estimated from such reports. This paper asks what features of the shape of a latent income distribution are identified from rounded reports, and how inference about them should be calibrated.

\paragraph{Model and contributions.}
Let $X$ be latent income and let the report be $Y=r_L(X)$, where $r_\ell$ rounds to the nearest multiple of $\ell$ and the level $L$ is drawn from a finite ladder $\lad\ni 1$ with probabilities $\pi_\ell(x)=\Prob(L=\ell\mid X=x)$ that may depend on income. Write $W=\lcm\lad$.

\emph{Identification of the mechanism.} Conditionally on $L=\ell$ the report is $\ell$ times an integer, and all the arithmetic information it carries about the level is contained in $G=\gcd(Y,W)$. If the residue of that integer modulo $W/\ell$ is uniform, a condition on the income distribution that we bound through its characteristic function, then $\Prob(G=d\mid L=\ell)=\ell\,\varphi(W/d)/W$ for $\ell\mid d\mid W$, with $\varphi$ Euler's totient, and the marginal rounding probabilities are identified from the divisibility pattern of the reports (Proposition~\ref{prop:totient}). Conditionally on its pattern, the residue of a report modulo $W$ is uniform over the residues with that pattern, whatever the mechanism, which yields a specification test with an exact conditional calibration and no nuisance parameter (Theorem~\ref{thm:spec}), sensitive to the failure that matters: reports crowding into part of the period $[0,W)$ because income does not span many periods of $W$.

\emph{Resolution.} Rounding to a level that divides the cell width moves mass only across cell boundaries, and if boundary reports are split between the two cells the reported cell probabilities equal the latent ones up to a term of second order in the level (Lemma~\ref{lem:split}). At the least common multiple $w_0$ of the active levels the latent cell vector is therefore identified without any model of how respondents choose a level, and since binning preserves unimodality (Lemma~\ref{lem:binning}), unimodality of the latent density becomes the hypothesis that a multinomial probability vector lies in the cone of unimodal sequences.

\emph{Inference at that resolution.} We measure departure from the cone by the discrete excess mass and calibrate it by resampling from a projection onto the cone after adjacent cells whose difference is within $\kappa_n$ standard errors have been pooled, as in generalized moment selection \citep{AndrewsSoares2010}. This bootstrap reproduces the limit law of the statistic at every fixed point of the null, including points on its boundary, where resampling from the unmodified projection does not (Theorem~\ref{thm:boot}).

\emph{Finer resolution, and its price.} Below $w_0$ the reported cell vector is a known linear map of the latent one once the profile $\pi_\ell(\cdot)$ is known, and for ladders totally ordered by divisibility the map is invertible at every width down to the finest active level (Theorem~\ref{thm:invert}). The profile needed is conditional on latent income, not on the report; band fits of the mixture identify its band averages, but the inversion amplifies small errors in any estimated profile and in our simulations holds its level at large sample sizes only when the profile is known, so we present it as a way to use a profile measured outside the survey, gated by the specification test.

\emph{Evidence.} On every American earnings variable we examine the specification test rejects the mechanism, for an identifiable reason: the ladder's least common multiple is \$50{,}000 in most samples and between two fifths and two thirds of reports lie below it, so the residues cannot be close to uniform, and in the ACS the Census Bureau's own rounding of released amounts violates the condition by construction. The certified width then leaves four to ten cells: the reports support no statement about the modality of American earnings, and the procedure says so rather than returning a rejection that reflects the questionnaire. On monthly housing costs the ladders are fine; a valley at the bottom of the distribution of gross rent does not survive a coarser grid and is not retained. On disposable income of one-person households in the CPS the lattice enters with transfer-dependent households, and once they are removed the specification test accepts the mechanism in four of five years. A thinning between \$25{,}000 and \$32{,}000 is significant only among respondents whose gross earnings are multiples of \$5{,}000: it is rounding in gross earnings carried through the tax computation to amounts that are no longer round, which a specification test applied to the constructed variable cannot see; the test belongs on the amount that was reported.

\paragraph{Related literature.}
Demography summarizes digit preference in reported ages by indices of the share of reports ending in zero or five \citep{Whipple1919,Myers1940,AHearnBatenCrayen2009}. \citet{HeitjanRubin1990,HeitjanRubin1991} place coarsened data in the missing-data framework; a mechanism in which the level depends on income is not coarsening at random and must be estimated. \citet{WrightBray2003}, \citet{WangHeitjan2008} and \citet{BarLillard2012} fit rounding as a mixture over levels jointly with a parametric model for the latent variable; \citet{CamardaEilersGampe2008} estimate digit preference jointly with a penalized latent density in the composite link model of \citet{Eilers2007}; \citet{GrossRendtel2016}, \citet{ZinnWurbach2016} and \citet{DrechslerKiesl2016} treat heaped income by stochastic EM or multiple imputation, again jointly with a model of the density. These joint approaches replace our residue condition by smoothness of the latent density at the scale of the finest level; neither assumption dominates, and Section~\ref{sec:data} shows data in which ours fails and is detected. The closest antecedents of our identification argument are \citet{BallTorousTschoegl1985} and \citet{Harris1991}, who estimate the probabilities with which prices are quoted to successively coarser resolutions from the clustering of observed prices; we make the divisibility structure explicit, its law being a matrix of totients whose triangular structure, with the positive definiteness of greatest-common-divisor matrices \citep{Smith1875,BeslinLigh1989}, gives identification, and whose ancillary part supplies a specification test. The Fourier analysis of rounding goes back to \citet{Sheppard1898}; \citet{Janson2006} treats it in generality, and \citet{Heitjan1989} reviews inference from grouped data. In economics, rounding has been read as a signal of uncertainty \citep{ManskiMolinari2010,KleinjansVanSoest2014,RuudSchunkWinter2014,Binder2017}, the interval a report represents being the object of interest.

The modality literature \citep{Silverman1981,HartiganHartigan1985,MullerSawitzki1991,ChengHall1998,HallYork2001,FisherMarron2001,DumbgenWalther2008,AmeijeirasAlonsoCRC2019} is built for continuous data. At the certified width our problem is order-restricted inference for a multinomial \citep{RobertsonWrightDykstra1988,SilvapulleSen2005}; maximum likelihood estimation of a unimodal probability mass function is studied by \citet{BalabdaouiJankowski2016}. The failure of the bootstrap on the boundary of a parameter space is due to \citet{Andrews2000}, and our calibration follows the moment-selection approach of \citet{AndrewsSoares2010} and the directional-differentiability framework of \citet{FangSantos2019}. Partial identification from interval-valued reports is developed by \citet{ManskiTamer2002}.

Proofs not given in the text, and additional tables, are in the supplementary material.
